\DocumentMetadata{
  pdfversion=2.0,
  pdfstandard=ua-2,
  lang=en-US,
  tagging=on,
  tagging-setup={math/setup=mathml-SE,tabsorder=structure}
}
\documentclass[11pt,letterpaper]{article}
\usepackage[margin=0.68in,includefoot]{geometry}
\usepackage{newcomputermodern}
\usepackage{amsmath,amscd,mathtools}
\usepackage{tikz,adjustbox}
\providecommand{\square}{\mdlgwhtsquare}
\usepackage[hidelinks]{hyperref}
\hypersetup{
  pdftitle={Algebra First-Year Exam, May 2015, Part 2},
  pdfauthor={University of Florida Department of Mathematics},
  pdfsubject={Graduate mathematics examination},
  pdfdisplaydoctitle=true,
  bookmarksopen=true
}
\setcounter{secnumdepth}{0}
\setlength{\parindent}{0pt}
\setlength{\parskip}{0.28em}
\setlength{\emergencystretch}{3em}
\setlength{\leftmargini}{2.15em}
\setlength{\leftmarginii}{2.65em}
\setlength{\leftmarginiii}{3em}
\pagestyle{empty}
\raggedbottom
\allowdisplaybreaks
\NewTaggingSocketPlug{block/list/label}{examproblem}{%
  \tagstructbegin{tag=Lbl}\tagstructend%
  \tagstructbegin{tag=\LBody}%
  \tagstructbegin{tag=\ProblemHeadingTag,title-o={\ProblemHeadingText}}%
  \tagmcbegin{tag=\ProblemHeadingTag,actualtext={\ProblemHeadingText}}%
  #2%
  \tagmcend%
  \tagstructend%
}
\newenvironment{examproblems}[2]{%
  \def\ProblemHeadingTag{#2}%
  \AssignTaggingSocketPlug{block/list/label}{examproblem}%
  \begin{enumerate}%
  \setcounter{enumi}{#1}%
  \renewcommand{\labelenumi}{\textbf{\theenumi.}}%
  \setlength{\topsep}{0.35em}%
  \setlength{\partopsep}{0pt}%
  \setlength{\itemsep}{0.48em}%
  \setlength{\parsep}{0.18em}%
}{\end{enumerate}}
\newenvironment{examsubparts}{%
  \AssignTaggingSocketPlug{block/list/label}{default}%
  \begin{enumerate}%
  \setlength{\topsep}{0.18em}%
  \setlength{\partopsep}{0pt}%
  \setlength{\itemsep}{0.16em}%
  \setlength{\parsep}{0.1em}%
}{\end{enumerate}}
\newenvironment{examinstructions}{%
  \par\begingroup\itshape
}{%
  \par\endgroup\vspace{0.18em}
}
\makeatletter
\renewcommand\subsection{\@startsection{subsection}{2}{0pt}%
  {-1.25ex plus -.25ex minus -.1ex}%
  {0.55ex plus .1ex}%
  {\normalfont\large\bfseries}}
\renewcommand{\@maketitle}{%
  \newpage\null\vskip -1em%
  {\footnotesize\bfseries
    \href{https://www.ufl.edu/}{University of Florida}, \href{https://math.ufl.edu/}{Department of Mathematics}\par}%
  \vskip -0.5em%
  \tagstructbegin{tag=H1,title={Algebra First-Year Exam, May 2015, Part 2}}%
  \tagpdfparaOff%
  \tagmcbegin{tag=H1}%
    {\LARGE\bfseries\href{https://gma.math.ufl.edu/exams/algebra-first-year/}{Algebra First-Year Exam}, May 2015, Part 2\par}%
  \tagmcend%
  \tagpdfparaOn%
  \tagstructend%
  \vskip 0.25em%
  \tagmcbegin{artifact}%
    \hrule height 1pt%
  \tagmcend%
  \vskip 1em%
}
\makeatother
\title{Algebra First-Year Exam, May 2015, Part 2}
\author{}
\date{}
\begin{document}
\maketitle
\thispagestyle{empty}
\begin{examinstructions}
Answer four problems. If you turn in more than four, only the first four will be graded. Within reason, you may use theorems as long as you state them clearly.
\end{examinstructions}
\begin{examproblems}{0}{H2}
\def\ProblemHeadingText{Problem 1.}
\item
This problem has three parts.
\begin{examsubparts}
\item[\textnormal{(i)}]
Find all units in the ring \(R=\mathbb{Z}[\sqrt{-7}]\).
\item[\textnormal{(ii)}]
Show that~\(R\) is not a unique factorization domain.
\item[\textnormal{(iii)}]
Show that \((2,1+\sqrt{-7})\) is a maximal ideal of~\(R\).
\end{examsubparts}
\def\ProblemHeadingText{Problem 2.}
\item
Suppose~\(R\) is a commutative ring and \(I\subset R\) is an ideal. Show that the set of \(r\in R\) such that \(r^n\in I\) for some \(n>0\) is an ideal of~\(R\).
\def\ProblemHeadingText{Problem 3.}
\item
Let~\(R\) be an integral domain and~\(M\) a finitely generated~\(R\)-module. We say that~\(M\) has rank~\(n\) if there are~\(n\) elements \(m_1,\ldots,m_n\in M\) that are linearly independent, while any subset of~\(M\) with cardinality greater than~\(n\) is linearly dependent. Show that~\(M\) has rank~\(n\) if and only if~\(M\) has a submodule \(N\subseteq M\) such that (1)~\(N\) is free of rank~\(n\), and (2) \(M/N\) is torsion.
\def\ProblemHeadingText{Problem 4.}
\item
Which of the following polynomials are irreducible? Explain.
\begin{examsubparts}
\item[\textnormal{(i)}]
\(X^3+5X+1\in\mathbb{Q}[X]\)
\item[\textnormal{(ii)}]
\(X^4+X+2\in\mathbb{Q}[X]\)
\item[\textnormal{(iii)}]
\(X^5+Y^5-1\in\mathbb{C}[X,Y]\)
\end{examsubparts}
\def\ProblemHeadingText{Problem 5.}
\item
Suppose~\(M\) is an \(n\times n\) matrix with coefficients in an algebraically closed field. Show that the following statements are equivalent:
\begin{examsubparts}
\item[\textnormal{1.}]
\(M\) is nilpotent, i.e. \(M^k=0\) for some \(k\geq 0\).
\item[\textnormal{2.}]
\(0\) is the only eigenvalue of~\(M\).
\item[\textnormal{3.}]
\(M\) is similar to an upper triangular matrix with 0s on the diagonal.
\end{examsubparts}
\def\ProblemHeadingText{Problem 6.}
\item
Find representatives for each similarity class of matrices in \(M_5(\mathbb{C})\) whose minimal polynomial is \(X^3+X^2-X-1=(X+1)^2(X-1)\).
\end{examproblems}
\end{document}
